\chapter{Teraquop volume}\label{ch:teraquop-volume}

The key metric we use to evaluate performance throughout this work is the \textit{teraquop volume}.
It is the natural generalisation of the \emph{teraquop footprint} \cite{gidney2021fault} from a `space only' picture to a `full spacetime' picture.

\begin{definition}
    Consider a family (set) of circuits whose $j$-th member is defined on $n_j$ qubits and runs for an amount of time $T_j$.
    Given a noise model and a decoder,
    the \textbf{teraquop volume} is the smallest volume $n_j T_j$ such that the probability of any spacelike and timelike logical error in the $j$-th circuit is $10^{-12}$.
\end{definition}

For a given circuit, noise model and decoder, the \emph{overall logical error rate} is the probability of any logical error -- whether spacelike or timelike -- occurring after the decoder has been run and its suggested corrections applied.
In practice, an approximation of this value will be found numerically.
If there is no $(n_j, T_j)$ such that the overall logical error rate is less than $10^{-12}$,
then we say the teraquop volume is infinite.

In the rest of this work we describe a DCCC with $n_M$ rows and $n_E$ columns of hexagons implemented for $h$ rounds of measurements as an \textit{$(n_M, n_E, h)$-block} of this DCCC.
(These rows are admittedly quite wonky -- for reference, in \Cref{fig:edge_and_face_operator} we showed a DCCC with $n_M = 3$ rows of and $n_E = 4$ columns.)
We say it has \textit{block height} $h$, \textit{block width} $n_E$ and \textit{block depth} $n_M$.
Then we'll write $\tilde{n}_M$, $\tilde{n}_E$, $\tilde{h}_M$, and $\tilde{h}_E$ for the minimum values of $n_M$, $n_E$, $h_M$, and $h_E$ such that the corresponding logical error rates are below $10^{-12}$.
An $(\tilde{n}_M, \tilde{n}_E, \tilde{h})$-block will be referred to as a \textit{teraquop block},
with \textit{teraquop height} $\tilde{h} = (\tilde{h}_E + \tilde{h}_M)/2$, and so on.